% ============================================================
% 04-ket-luan.tex — Kết luận và hướng phát triển (Bình)
% Số liệu dùng macro \soLieu... (results/so-lieu*.tex). Bốn câu hỏi CH1–CH4 khớp Mở đầu và Tiểu kết Chương 1.
% ============================================================
\chapter*{KẾT LUẬN VÀ HƯỚNG PHÁT TRIỂN}
\addcontentsline{toc}{chapter}{KẾT LUẬN VÀ HƯỚNG PHÁT TRIỂN}

\section*{Kết quả đạt được}

Tiểu luận đi theo 2 giai đoạn. Nhóm bắt đầu bằng cách ghép phổ biến, dùng PSO tìm tâm khởi tạo cho FCM, và thấy cách này
không cải thiện FCM trên dữ liệu tín dụng, đúng như các so sánh độc lập trong tài liệu. Nhóm chuyển sang dùng PSO tìm trọng số
thuộc tính cho FCM, đề xuất 3 lựa chọn thiết kế cho bài toán này, so 5 phiên bản PSO và kiểm trọng số trên mẫu độc lập.
Câu trả lời cho 4 câu hỏi nghiên cứu, với số liệu ở Bảng~\ref{tab:ch3-tong-ket}, như sau.

\emph{CH1 (hướng ban đầu).} PSO tìm tâm khởi tạo không giúp FCM: các bản cải tiến tối ưu tốt hơn PSO gốc, nhưng sau khi FCM
chạy, mọi phiên bản cho cùng nghiệm với FCM khởi tạo k-means++ mà chậm hơn (khoảng 1,7 lần trên Việt Nam, 24 lần trên Taiwan Default). Nguyên nhân là hàm mục tiêu trên dữ liệu
đã thử hầu như chỉ có 1 cực tiểu, nên PSO và FCM làm giảm cùng một hàm. Các bài báo PSO--FCM báo kết quả tốt hơn chủ yếu vì
so với FCM khởi tạo ngẫu nhiên và không cùng ngân sách.

\emph{CH2 (đề xuất).} PSO tìm trọng số làm các phân khúc tách rủi ro tốt hơn FCM thường có ý nghĩa trên 2 bộ chính (Việt Nam
ở mọi số cụm, Taiwan Default với 3 và 4 cụm); với 2 cụm trên Việt Nam, chênh tỉ lệ vỡ nợ giữa các cụm tăng từ
\soLieuETamVietNamCHaiFcmKppChenhTiLe{} lên \soLieuETamVietNamCHaiPsoTrongSoChenhTiLe. So với tìm ngẫu nhiên cùng ngân sách, PSO
tối ưu fitness tốt hơn ở đa số cấu hình, nhưng chỉ tách rủi ro tốt hơn ở vài cấu hình, chủ yếu trên Việt Nam; trọng số tính bằng công thức không thay được PSO trên Taiwan Default.
Lợi ích này đo so với FCM $m = 2$: với ba, 4 cụm, K-means không trọng số vẫn tách rủi ro tốt hơn trên Việt Nam (và theo
chênh tỉ lệ vỡ nợ trên Taiwan Default với 4 cụm), nên FCM có trọng số chỉ đáng dùng khi cần độ thuộc. Trọng số tìm
được có nghĩa kinh tế: lịch sử trả nợ trên Taiwan Default, định giá và đòn bẩy trên Việt Nam.

\emph{CH3 (so sánh các phiên bản).} Cả 5 phiên bản đều hơn tìm ngẫu nhiên cùng ngân sách, nhưng không bản cải tiến nào hơn
PSO gốc, ngược với vai trò tìm tâm. Trong vai trò tìm trọng số, không phát hiện khác biệt giữa các phiên bản sau hiệu chỉnh Holm, nên
chọn phiên bản nào không phải yếu tố quyết định.

\emph{CH4 (kiểm chứng và giới hạn).} Trọng số học trên một mẫu vẫn cải thiện phân khúc trên mẫu kia so với FCM thường; theo chênh tỉ lệ
vỡ nợ, kết quả không khác học ngay trên mẫu đó ở \soLieuEMuoiChenhTiLePsoTaiChoBang{} trên \soLieuEMuoiSoCauHinh{} cấu hình, còn theo AUC
kém hơn ở \soLieuEMuoiAucCungPsoTaiChoKem{} cấu hình. Lợi ích không phổ quát: trên các bộ mà FCM thường đã gần sụp và với 2 cụm (kể cả Taiwan Default, mức chênh
nhỏ), PSO không có ích hoặc kém hơn, vì fitness không dùng nhãn.

Về thực hành, khi phân khúc khách hàng tín dụng bằng FCM, không nên dùng PSO tìm tâm: khởi tạo k-means++ là đủ, và nên chẩn
đoán số cực tiểu trước khi nghĩ tới PSO cho tâm. Nên dùng PSO để tìm trọng số thuộc tính theo đề xuất, và chỉ áp trọng số
khi nó cải thiện phân khúc trên một mẫu có nhãn lịch sử không dùng để tìm trọng số (mục~\ref{sec:de-xuat}).

\section*{Hạn chế}

\begin{itemize}
  \item Phần trọng số dùng một fitness (Silhouette trên không gian gốc, phạt cụm dưới 5\%), một số mũ $m = 2$ và $c$ từ 2 đến
  4; với $m = 2$ trên Taiwan Default, độ thuộc vẫn chia gần đều ngay cả khi phân hoạch cứng đã tách rủi ro tốt hơn.
  \item Hai chỉ số chính dùng nhãn để đánh giá và tính trên chính mẫu được phân cụm; TN5 kiểm trên mẫu độc lập nhưng chỉ cho
  2 bộ chính. Bộ Việt Nam gồm các quan sát doanh nghiệp--năm không độc lập; 2 nửa được chia theo doanh nghiệp để giảm ảnh
  hưởng này.
  \item Đối chứng chính là FCM với $m = 2$, vốn gần suy biến khi có từ 3 cụm; K-means không trọng số được thêm làm mốc tham
  chiếu và tách rủi ro tốt hơn ở các cấu hình này. Chưa thử PSO tìm trọng số cho K-means hay cho FCM với $m$ nhỏ hơn.
  \item Kiểm định lấy seed làm đơn vị, nên chỉ đo dao động của thuật toán trên mẫu đã cho, không đo dao động do lấy mẫu khách
  hàng. Nhãn của bộ Việt Nam được định nghĩa từ chính các đại lượng tài chính có trong thuộc tính (vốn chủ sở hữu, EBITDA, lãi
  vay) và gán theo doanh nghiệp, nên mức tách rủi ro trên bộ này có thể cao hơn khi dự báo thật.
  \item Chỉ có 2 bộ dữ liệu chính, trong đó bộ Việt Nam và cả 2 mẫu Taiwan Default đã được dùng ở các thí nghiệm thử; 5 bộ phụ cho thấy kết quả không khái quát
  được cho mọi dữ liệu tín dụng.
  \item Silhouette của fitness tính trên mẫu 1.000 điểm nên có nhiễu mẫu; ngân sách 500 và 1.000 lần tính fitness nhỏ so với
  các bài đã đọc, và các phiên bản dùng chung tham số thay vì tham số riêng của từng bài.
  \item Phần tìm tâm dùng $m = 2$, $c = 3$ trong thí nghiệm chính, một chế độ gần suy biến; kiểm lại với $m = 1{,}5$ cho cùng
  kết luận nhưng chỉ trên 2 bộ dữ liệu. ``Số cực tiểu'' là số mức $J_m$ quan sát được, không phải số cực tiểu thật.
  \item Thời gian được đo khi nhiều tiến trình chạy song song và máy còn chạy việc khác, nên chỉ chênh lệch nhiều lần là đáng
  tin.
  \item Danh sách thuộc tính gần trùng của bộ Việt Nam được chọn thủ công và bỏ sót X10 (tương quan 0,966 với X8).
\end{itemize}

\section*{Hướng phát triển}

Ba hướng tiếp nối trực tiếp phần trọng số: tìm trọng số cùng lúc với $m$ và $c$ thay vì cố định chúng, thử các fitness khác
Silhouette (chẳng hạn chỉ số mờ như XB khi các tâm tách rõ), và cho phép trọng số riêng cho từng cụm như WFCM của bài
06~\cite{liu2021msgwo}. Hai hướng khác là xử lý dữ liệu hỗn hợp số và định danh để FCM không sụp trên các bộ tín dụng cá nhân,
và dùng các phân khúc có trọng số làm đầu vào cho mô hình chấm điểm theo cụm như Bosker và cộng sự~\cite{bosker2025clustered}.
